\paragraph{RQ3: interpretation and limitations.}
Our categories do not consistently improve on Shin, and categories plus joint angles do not fully replace his directions under the tested settings. The small world-angle improvement depends on the learner and coordinate source. Permutation dependence should not be read as the gain from adding a feature.

The short left/right and G/H explorations illustrate limits of synthetic changes: improved artificial scores do not establish better live recognition. They are practical follow-ups, not additional independent validation of the main feature study. Rigid G/H rotation preserves geometry and does not recreate new detector errors under occlusion.

The folds were inspected during development and are not verified people, so these are exploratory results rather than an untouched final test. Detection recovery, accepted-image classification and selected webcam errors have different denominators and should be interpreted separately.
